\begin{annexitem}{annex:p3.n-a-006.2.uniform}{thm:p03-bl-uniform}{Compactness supplies the uniform bound}
\begin{proof}
Suppose no bound exists. For every \(q\) there is a parameter tuple,
in a model of the complete theory of \((\mathcal U,\bar c)\), for which
both the solution set and its complement within \(D\) have more than
\(q\) elements. Name a new tuple \(\bar b\) and require, for every
positive integer \(r\), at least \(r\) distinct solutions and at least
\(r\) distinct nonsolutions. Each finite collection of requirements is
satisfiable, by taking \(q\) larger than all its bounds. Compactness
gives a model with an instance having infinitely many of both. By the
elementary diagram, the same argument can be made in an elementary
extension of \((\mathcal U,\bar c)\): each finite existence sentence true
in a model of its complete theory is true in \((\mathcal U,\bar c)\).
This contradicts strong minimality. Increase the resulting bound if
necessary so it is nonnegative. A fibre with more than \(q\) elements
cannot be finite, so is cofinite; a cofinite fibre in an infinite \(D\)
has more than \(q\) elements. The assertion of at least \(q+1\) distinct
solutions is first-order. Its negation expresses the finite alternative.
\end{proof}
\end{annexitem}

\begin{annexitem}{annex:p3.n-a-006.2.pregeometry}{thm:p03-bl-pregeometry}{Finite fibres and exchange}
\begin{proof}
Extensivity follows from the singleton formula \(x=a\), monotonicity
from inclusion of available parameters, and finite character from the
finiteness of a formula's parameter tuple. Algebraic closure is
transitive: suppose \(a\) satisfies a finite-fibre formula over
\(b_1,\ldots,b_s\), each \(b_i\) algebraic over \(A\). Choose finite
\(A\)-definable sets containing the \(b_i\), and conjoin to the formula
the assertion that its fibre has at most its actual finite number of
solutions. Existentially quantify the \(b_i\) over those finite sets.
The resulting \(A\)-formula contains \(a\) and defines a subset of a
finite union of uniformly finite fibres, hence a finite set. Thus
\(\operatorname{acl}(\operatorname{acl}(A))=\operatorname{acl}(A)\).
Intersecting with \(D\), and retaining \(\bar c\) in the base, gives
idempotence of the stated closure.

For exchange, write \(C\) for the full fixed base together with \(X\),
and assume \(a\in\operatorname{acl}(C b)\setminus\operatorname{acl}(C)\).
Choose a finite tuple from \(C\) and a formula \(\rho(x,y)\), with those
parameters, whose \(x\)-fibre at \(b\) contains \(a\) and has at most
\(r\) solutions, where \(r\geq1\). Replace \(\rho\) by the conjunction
of \(x,y\in D\), the old formula, and the first-order assertion that
its \(x\)-fibre has at most \(r\) elements. The new relation \(R\)
still contains \((a,b)\), and every \(x\)-fibre of \(R\) has at most
\(r\) elements.

Suppose \(b\notin\operatorname{acl}(C a)\). The \(y\)-fibre
\(R(a,D)\) contains this nonalgebraic element, so it is not finite and
is therefore cofinite. By Lemma~\ref{thm:p03-bl-uniform}, the condition
that \(R(x,D)\) is cofinite is definable over the finite parameters
from \(C\). Let \(H\subseteq D\) be the set it defines. Since
\(a\in H\) but \(a\notin\operatorname{acl}(C)\), \(H\) cannot be
finite; it is cofinite. Choose \(r+1\) distinct elements of \(H\).
Their \(y\)-fibres are cofinite, so their intersection is cofinite and
in particular nonempty. A \(y\) in that intersection has at least
\(r+1\) distinct points in its \(x\)-fibre, contradicting the
bound \(r\). Hence \(b\in\operatorname{acl}(C a)\). All definitions
used only finitely many parameters at a time, so the proof covers an
arbitrary additional fixed base, not just a finite \(C\).
\end{proof}
\end{annexitem}

\begin{annexitem}{annex:p3.n-a-006.2.dimension}{thm:p03-bl-dimension}{Replacement and equality of basis cardinalities}
\begin{proof}
The union of an inclusion-chain of independent sets is independent: any
algebraic dependence uses finitely many elements and would already occur
in a member of the chain. Zorn's lemma therefore extends any independent
set to a maximal one \(I\). If \(a\notin\operatorname{cl}_{\bar c}(I)\),
then \(I\cup\{a\}\) is independent. Indeed, a dependence of an old
\(i\) on the other old elements together with \(a\), when there was
no old dependence, would by exchange make \(a\) depend on \(I\).
This contradicts the choice of \(a\). Maximality thus implies that
\(I\) spans \(D\).

Here is the finite replacement argument. Let an independent list
\(a_1,\ldots,a_s\) lie in the closure of a finite set \(J\) with
\(r\) elements. Maintain a spanning set \(J_i\) of size \(r\)
containing \(a_1,\ldots,a_i\), beginning with \(J_0=J\).
If \(a_{i+1}\) is already in \(J_i\), no replacement is needed.
Otherwise choose a minimal subset \(F\subseteq J_i\) spanning
\(a_{i+1}\). It is not contained in \(\{a_1,\ldots,a_i\}\), since
that would contradict independence. Choose
\(j\in F\setminus\{a_1,\ldots,a_i\}\). Minimality and exchange give
\[
 j\in\operatorname{cl}_{\bar c}
       ((F\setminus\{j\})\cup\{a_{i+1}\}).
\]
Thus \(J_{i+1}=(J_i\setminus\{j\})\cup\{a_{i+1}\}\) still spans
\(J_i\), and hence its previous closure, and has size \(r\).
The process cannot place more than \(r\) distinct independent elements
in a set of size \(r\); hence \(s\leq r\). This also covers
\(r=0\), when there are no nonalgebraic independent elements.

For two bases \(I,J\), every finite subset of \(I\) lies in the
closure of a finite subset of \(J\), by finite character. If \(J\)
is finite, the preceding inequality shows \(\abs{I}\leq\abs{J}\).
If \(J\) is infinite, partition the membership calculation as
\[
 I=\bigcup_{F\subseteq J\text{ finite}}
        (I\cap\operatorname{cl}_{\bar c}(F)).
\]
Each intersection is an independent set of size at most \(\abs{F}\),
so finite. There are \(\abs{J}\) finite subsets of an infinite \(J\),
and therefore \(\abs{I}\leq\abs{J}\). Reversing \(I,J\) gives
equality, including the cases of finite and empty bases.

In a countable language, the formulas over \(\bar c\cup I\) number
at most \(\max(\aleph_0,\abs{I})\). Each algebraic element belongs
to the finite solution set of one of them. Their union bounds the
cardinality of the closure by the same infinite cardinal. Since
\(I\subseteq D=\operatorname{cl}_{\bar c}(I)\), uncountability of
\(D\) forces \(\abs{I}=\abs{D}\).
\end{proof}
\end{annexitem}

\begin{annexitem}{annex:p3.n-a-006.2.independent-types}{thm:p03-bl-independent-types}{Generic decisions and finite-tuple induction}
\begin{proof}
A \(B\)-definable subset of \(D\) is finite or cofinite. An element
outside \(\operatorname{acl}(B)\) belongs to none of the finite
ones, and to all of the cofinite ones because their complements are
finite. These decisions determine its complete type over \(B\).

Induct on the ordered tuple length. For the step, suppose the two
initial tuples have the same type over \(B\). For any formula
\(\theta(x,\bar z,\bar b)\), with a finite \(\bar b\) from \(B\),
the property that \(\theta(D,\bar z,\bar b)\) is cofinite is
first-order in \(\bar z\) by the uniform-bound lemma. Thus it has
the same truth value at the two initial tuples. The next elements
are nonalgebraic over their respective initial tuples and \(B\),
so both satisfy \(\theta\) exactly in the cofinite case. They
therefore extend the tuples to the same complete type. Full
independence implies this sequential nonalgebraicity by monotonicity.
Conversely, sequential nonalgebraicity implies full independence.
If there were a dependence, choose an inclusion-minimal finite
dependent set \(C\) of coordinates. Some \(c\in C\) belongs to the
closure of \(C\setminus\{c\}\), and every proper subset of \(C\)
is independent. For any \(d\in C\setminus\{c\}\), minimality says
that \(c\) is not in the closure of \(C\setminus\{c,d\}\).
Exchange over this smaller set therefore puts \(d\) in the closure
of \(C\setminus\{d\}\). Thus every member of \(C\), including
its last coordinate in the original order, depends on the others.
That last coordinate depends only on earlier coordinates,
contradicting its sequential condition. Finally, every formula in
the domain of an independent-set bijection uses finitely many
elements. Apply the finite-tuple conclusion to those elements and
to its finitely many parameters from \(B\).
\end{proof}
\end{annexitem}

\begin{annexitem}{annex:p3.n-a-006.2.whole-models}{thm:p03-bl-whole-models}{Algebraic extension of base maps and the countable spectrum}
\begin{proof}
For a model \(M\) and a basis \(I\), the pregeometry result gives
\(M=\operatorname{acl}(I)\). Every algebraic solution over parameters
in a model already lies in that model: if a formula has exactly
\(r\) solutions in an elementary extension, the first-order sentence
asserting exactly \(r\) solutions supplies all \(r\) in the model.
Consequently closure computations here agree with those in a common
elementary extension.

We first give the extension step without a saturation assumption.
Let \(f:A\to N\) be a partial elementary map whose domain contains
\(I\), and let \(a\in M\setminus A\). A formula over a finite
tuple from \(I\) defines a finite set containing \(a\). Its image
defines a finite set \(F\) in \(N\) of the same size. Each finite
collection of formulas in the image of \(\operatorname{tp}(a/A)\)
has a realization in \(F\): the corresponding existential formula
over finitely many elements of \(A\) is true in \(M\), so in \(N\)
by elementarity of \(f\). The resulting family of subsets of finite
\(F\) has the finite-intersection property; its total intersection
is nonempty. Choose a member as \(f(a)\). The extended map is
elementary. Its type includes inequalities from every old domain
element, so the new value is not in the old image.

To compare different models, compactness places elementary copies
of them in a common model of \(T\). Use disjoint new names for their
elementary diagrams. Each finite collection from either diagram,
after existentially quantifying its names, is a sentence true in
every model of complete \(T\). The two finite collections can
therefore be realized together, so the combined diagrams are
consistent. The generic-type proof can now be applied to both
copies with the empty base.
A bijection of bases is elementary by
Lemma~\ref{thm:p03-bl-independent-types}. Extend it alternately in
the two directions using the preceding finite-set step, through
well-orderings of the universes; at a limit take the union. Each
formula mentions only finitely many mapped elements, so the union
is elementary. The result is onto and is an isomorphism.
An isomorphism preserves algebraic closure and independence, so the
converse follows from equality of basis cardinalities.

The same extension construction in just the forward direction
extends any elementary basis map into any model, proving primeness
over \(I\). If \(M_0\preccurlyeq M\) contains \(I\), it contains
all the finite algebraic solution sets over \(I\), so
\(M_0=M\). This proves minimality.

For the existence calculation, every model has elementary extensions
with arbitrarily long independent lists. At a successor step the
partial type excluding every finite definable subset over the
previous parameters is finitely satisfiable: a finite union of
finite subsets cannot exhaust the infinite universe. Compactness
with the elementary diagram realizes it in an elementary extension.
Continue along an ordinal and take elementary unions at limits.
The resulting list is independent by exchange and the sequential
criterion used in the previous proof.

If \(A\) is an infinite algebraically closed subset of a model of
\(T\), it is an elementary submodel. It is a substructure because
values of function symbols on tuples from \(A\) are definable, hence
algebraic, over those tuples. For any formula with parameters in
\(A\) and with a witness, its solution set is finite or cofinite.
In the finite case all solutions lie in \(A\); in the cofinite
case some witness lies in infinite \(A\). The Tarski--Vaught test
now proves elementarity. Hence the algebraic closure of any infinite
independent set is a model having that set as a basis.

In a countable language such closures of infinite bases of cardinal
\(\kappa\) have cardinal \(\kappa\), by the closure cardinal bound.
Every uncountable model has dimension equal to its cardinality,
and the basis-isomorphism result proves categoricity there. A
countable model has a finite or countably infinite basis. If no
finite-dimensional model exists, the latter case is its only
isomorphism type. If one exists, let \(n_0\) be the least finite
dimension of a model. Extend a basis of an \(n_0\)-dimensional
model by \(r\) successive independent elements. The closure of the
enlarged basis contains the old infinite model, so is infinite and
is elementary by the preceding test. It has dimension \(n_0+r\).
There is also a model with a countably infinite basis. Minimality
of \(n_0\) excludes lower finite dimensions, and the isomorphism
classification supplies exactly one type at each listed dimension.
This proves the dichotomy for this whole-universe hypothesis.
\end{proof}
\end{annexitem}

\begin{annexitem}{annex:p3.n-a-006.2.acf}{thm:p03-bl-acf-geometry}{One-variable field formulas and transcendence degree}
\begin{proof}
The quantifier-elimination theorem explicitly used here is
Theorem~\ref{thm:algebraically-closed-fields}. Every one-variable ring formula with
parameters is therefore equivalent to a finite Boolean combination
of polynomial equations. A polynomial in one variable is either
identically zero or has at most its degree many roots. Its zero set
is thus the entire field or finite. Boolean combinations of finite
and cofinite sets are finite or cofinite. The same argument holds
in every algebraically closed elementary extension, establishing
strong minimality. Such a field is infinite: a finite field with
elements \(a_1,\ldots,a_r\) would have no root for
\(\prod_{i=1}^r(x-a_i)-1\). Quantifier elimination also reduces a
sentence to a Boolean combination of equations between closed ring
terms. Their values lie in the prime field and their equality
depends only on the fixed characteristic. Thus all models agree on
sentences, proving completeness.

Let \(F\) be the subfield generated by the prime field and \(B\).
If \(a\) is field-algebraic over \(F\), clear denominators in its
nonzero polynomial equation. Its coefficients are rational
expressions in finitely many elements of \(B\); the nonzero
denominators are specified, and the equation defines finitely many
roots. Thus \(a\in\operatorname{acl}(B)\). Conversely, if a
formula over \(B\) defines a finite set containing \(a\), eliminate
quantifiers. Collect the finitely many nonzero polynomials occurring
in the resulting Boolean combination. Outside the union of their
root sets, every atom has a fixed truth value, so the combination
has a fixed truth value. That value must be false, since otherwise
the definable set would be cofinite and infinite. Therefore \(a\)
is a root of one of these nonzero polynomials, whose coefficients
lie in \(F\), and is field-algebraic over \(F\). If the list were
empty, the formula would be constant and could not define a finite
nonempty set. This establishes equality of the two closures.

Independence in this closure is exactly algebraic independence over
the prime field, and bases are transcendence bases. For each
\(n<\omega\), take an algebraic closure of the rational function
field in \(n\) independent variables over the prime field. Also
take an algebraic closure in countably infinitely many independent
variables. Each field is countable: there are countably many
rational expressions and polynomials, each polynomial has finitely
many roots, and the algebraic extension has a countable union of
finite root sets. In positive characteristic, even degree zero is
infinite, since a finite field cannot be algebraically closed.
These constructions give dimensions \(0,1,2,\ldots,\aleph_0\).
The preceding classification proves uniqueness at each dimension,
excludes any other countable case, and gives uniqueness in every
uncountable cardinal. Existence in the uncountable cardinals follows
in the same way from algebraically closing a rational function field
on that many variables; the cardinal bound is attained because
the independent variables remain distinct.
The field-theoretic existence of an algebraic closure is the
standard algebraic-closure theorem; this example uses it as an
external algebraic input \cite{Steinitz1910}, not as a consequence
of model-theoretic exchange.
\end{proof}
\end{annexitem}
